\section{Discussion}
\label{sec:discussion}

\paragraph{Prediction is not identification.}
The three layers of evidence agree. The data determine, at best, a combination of the two gut rates, and under
the area they determine it for few subjects (Sections~\ref{sec:timing} and~\ref{sec:more}); insulin
sensitivity is constrained from one side for most subjects and bounded for a minority, even when the model
is true (Section~\ref{sec:si}); and none of this changes how well the fitted model predicts held-out meals
(Section~\ref{sec:prediction}). A model whose $\ke$ and $\ka$ are held at their population values predicts
the held-out area as well as one whose values were fitted (H8), as Proposition~\ref{prop:moments} would lead
one to expect of an area. Predictive validation is therefore silent on whether a fitted parameter means
anything, and a twin whose parameters are shown to a person should carry an identifiability analysis with it.

\paragraph{What to report.}
The results suggest a short list for anyone fitting a mechanistic twin to a summary of free-living data.
\begin{enumerate}
\item Report profile intervals, not point estimates, and say whether each is bounded, one-sided or flat; a
  one-sided interval is the prior range speaking, so report it at more than one range.
\item Define the set of subjects over which a fraction is taken. Parameters that are not identified sit on
  bounds, so ``no parameter on a bound'' selects an empty set.
\item Report the identifiable combinations (here $\tau_1$ and $p$, under the trace) rather than the
  individual parameters they hide behind.
\item Run a replica in which the model is true. If the data do not bound a parameter there, they will not
  bound it in the real data for the same reason, and the interval is not evidence of misspecification.
\item Test the optimizer against a better-conditioned one; a fixed-step estimate that is not an optimum
  produces profile minima below the estimate, as we found for the coordinate fits (Section~\ref{sec:more}).
\item Treat clinical correlations of a fitted parameter with care: a correlation that survives only
  until the fasting level is partialled out is a statement about the fasting level.
\end{enumerate}

\paragraph{Design.}
Proposition~\ref{prop:moments} says what an observable must carry to identify the timing parameters: the
first moment for $\tau_1$ and the second for $p$. An area over a long window carries neither. A short window
carries some timing in the area but removes information about $\SI$; the trace carries both and is the only
observable in our ladder that bounds the combination for most subjects. The proposition also suggests, though
we did not test it, that an intervention which changes the gastric step without changing absorption (a liquid
in place of a solid meal, say) is the way to break the exchange symmetry by design rather than by assumption.
A design sweep over meal size and timing, which would quantify how much of the weakness a better protocol
could remove, was planned and not run.

\subsection{Limitations}
\label{sec:limits}

\begin{itemize}
\item \emph{One model family and one observable.} The analysis is of a Bergman-type model with a
  two-compartment gut, fitted to the iAUC and two richer observables of postprandial windows. The exchange
  symmetry is a property of any such cascade, but the second model class (Section~\ref{sec:dm}) was examined
  only for the iAUC and only on one cohort; models with other gut dynamics, other observables or other
  feedbacks may differ.
\item \emph{Cohorts.} The cohorts are of modest size (\resCohortCgmSubjects, \resCohortShanghaiSubjects{} and
  \resCohortHallSubjects{} subjects) and differ in sensor, window and how carbohydrate was obtained; prediction
  results are from two of them, and the replica and synthetic studies are built on CGMacros meal schedules.
  In the other two cohorts most estimates sit on a bound at the primary box, which limits how much the
  interior fractions can say.
\item \emph{The replica is an idealization.} Its truths are fitted values, which sit near the upper bound
  of the timing rates; its noise scale comes from the residuals of the real fit and so includes misfit; and its
  carbohydrate error is lognormal with a fixed coefficient of variation. The synthetic study with random
  truths removes the first limitation and not the others.
\item \emph{Noise model.} A single variance per subject and observable, frozen at a pilot fit, with
  first-order autocorrelation for the trace. We did not test heteroskedastic or non-Gaussian noise.
\item \emph{The criterion.} A bounded interval at a rise of $\delta$ is an asymptotic criterion; with a few
  tens of meals per subject and a hard box it is approximate, which is why the results are reported at
  three box widths and with a replica for calibration. Coverage of bounded intervals in the replica was
  below nominal.
\item \emph{Meals are treated as isolated.} Every meal starts with an empty gut. Residual gut content breaks
  the symmetry in proportion to its size (Proposition~\ref{prop:breaking}); we have not estimated how large it
  is in free-living data.
\item \emph{Post hoc hypotheses.} Hypotheses H13 to H22 were declared after earlier results were seen. Their
  thresholds were fixed before the corresponding analyses were run, but they were chosen knowing the
  direction of the earlier results.
\item \emph{Not run.} The sensitivity of the $\SI$ profile to heteroskedastic noise and a design sweep over
  meal size and timing were planned and not run (Appendix~\ref{app:notrun}).
\end{itemize}

\subsection{Conclusion}

A digital twin whose parameters are fitted to a postprandial summary of free-living CGM can predict
held-out meals as well as the data allow while leaving two of its three personalised parameters
unidentified and the third identified for a minority of people. The gut rates are separable only as a
combination and only from the trace; insulin sensitivity is only conditionally recoverable and is
sensitive to the carbohydrate log in proportion; and prediction, optimizer comparison and parameter
count are all blind to these facts. We report the analyses as an operational protocol and the failures as
findings, because the failures are the point: a good fit is not a measurement.
